\section{前向扩散过程}

\subsection{单步加噪}

\begin{definition}[前向扩散过程]
给定噪声调度 $\{\beta_1,\dots,\beta_T\}$，$\beta_t\in(0,1)$，前向过程按如下\emph{马尔可夫链}
逐步给数据加噪：
\begin{equation}
  q(\bx_t\mid\bx_{t-1})
  = \Ncal\Bigl(\bx_t;\ \sqrt{1-\beta_t}\,\bx_{t-1},\ \beta_t\bI\Bigr),
  \qquad t=1,2,\dots,T,
  \label{eq:forward-step}
\end{equation}
等价地写成重参数化的形式
\begin{equation}
  \bx_t = \sqrt{1-\beta_t}\,\bx_{t-1} + \sqrt{\beta_t}\,\beps_{t-1},
  \qquad \beps_{t-1}\sim\Ncal(\bzero,\bI)\ \text{相互独立}.
  \label{eq:forward-step-reparam}
\end{equation}
记 $\alpha_t = 1-\beta_t$，则整条链的联合分布为
$q(\bx_{1:T}\mid\bx_0)=\prod_{t=1}^{T}q(\bx_t\mid\bx_{t-1})$。
\end{definition}

\begin{remark}[为什么系数是 $\sqrt{1-\beta_t}$ 而不是 $1$]
式~\eqref{eq:forward-step-reparam} 里信号的系数与噪声的系数是一对平方和为 $1$ 的数：
$(\sqrt{1-\beta_t})^2+(\sqrt{\beta_t})^2=1$。若 $\bx_{t-1}$ 的各维独立、均值为 $0$、方差为 $1$，
那么 $\bx_t$ 的方差为 $(1-\beta_t)\cdot 1+\beta_t=1$，\emph{方差始终为 $1$}。
这个约定称为保方差（variance preserving）。反之若写成 $\bx_t=\bx_{t-1}+\sqrt{\beta_t}\beps$，
方差会随 $t$ 线性增长，到 $T=1000$ 时信号会被彻底淹没，网络也就学不动了。
\end{remark}

\subsection{任意时刻的闭式解}

逐层递推 $T$ 次显然不可行。但幸运的是，由于每一步都是高斯，整条链可以被\textbf{一步折叠}，
直接写出任意 $t$ 时刻 $\bx_t$ 关于 $\bx_0$ 的分布——这是扩散模型最核心的一条性质。

\begin{proposition}[前向过程的闭式解]
\label{prop:forward-closed}
对任意 $t=1,\dots,T$，
\begin{equation}
  q(\bx_t\mid\bx_0)=\Ncal\Bigl(\bx_t;\ \sqrt{\abt}\,\bx_0,\ (1-\abt)\bI\Bigr),
  \label{eq:forward-closed}
\end{equation}
即
\begin{equation}
  \boxed{\ \bx_t=\sqrt{\abt}\,\bx_0+\sqrt{1-\abt}\,\beps,\qquad \beps\sim\Ncal(\bzero,\bI).\ }
  \label{eq:forward-closed-reparam}
\end{equation}
\end{proposition}

\begin{pf}
对 $t$ 归纳。$t=1$ 时即为定义式~\eqref{eq:forward-step-reparam}。
设 $t-1$ 时命题成立，即 $\bx_{t-1}=\sqrt{\bar\alpha_{t-1}}\,\bx_0+\sqrt{1-\bar\alpha_{t-1}}\,\beps'$，
代入单步公式并合并同类项：
\begin{equation}
  \begin{aligned}
    \bx_t
    &= \sqrt{\alpha_t}\,\bx_{t-1}+\sqrt{\beta_t}\,\beps''
     = \sqrt{\alpha_t\bar\alpha_{t-1}}\,\bx_0
       + \underbrace{\sqrt{\alpha_t(1-\bar\alpha_{t-1})}\,\beps'+\sqrt{\beta_t}\,\beps''}_{\text{两个独立零均值高斯之和}} .
  \end{aligned}
  \label{eq:forward-induction}
\end{equation}
由附录~\ref{app:gauss-sum}，两个独立的零均值高斯之和仍是零均值高斯，其方差为两者方差之和：
\begin{equation}
  \alpha_t\,(1-\bar\alpha_{t-1})+\beta_t
  = \alpha_t-\alpha_t\bar\alpha_{t-1}+1-\alpha_t
  = 1-\alpha_t\bar\alpha_{t-1}
  = 1-\abt .
  \label{eq:forward-variance}
\end{equation}
又 $\sqrt{\alpha_t\bar\alpha_{t-1}}=\sqrt{\abt}$，故 $\bx_t$ 服从式~\eqref{eq:forward-closed}。
归纳完成。
\end{pf}

\begin{corollary}[起点的分布是已知的]
令 $t=T$ 得 $q(\bx_T\mid\bx_0)=\Ncal(\bx_T;\sqrt{\bar\alpha_T}\,\bx_0,(1-\bar\alpha_T)\bI)$。
只要调度使得 $\bar\alpha_T\to 0$，就有 $\bx_T\to\Ncal(\bzero,\bI)$ 且\emph{与 $\bx_0$ 无关}。
这保证了反向过程有一个已知的、可以直接采样的起点：从标准高斯里取一个噪声，就能开始生成。
\end{corollary}

\subsection{训练时的「想像中」与「实际上」}

按定义，要得到一个含噪样本 $\bx_t$ 似乎必须老老实实走 $t$ 步。但命题~\ref{prop:forward-closed} 告诉我们
可以\emph{一步算出}。这一点在教学材料中常被专门拿出来强调，因为它是扩散模型能够高效训练的
关键，也是初学者最容易忽略的地方\cite{lee-ddpm}。

\begin{keybox}[一句话记住 $\bx_t$ 的来源]
\begin{center}
  \textbf{想像中}：$\bx_0 \xrightarrow{\ \sqrt{\beta_1}\beps_1\ } \bx_1
                  \xrightarrow{\ \sqrt{\beta_2}\beps_2\ } \bx_2
                  \xrightarrow{\ \cdots\ } \bx_t$
                  \quad（逐步加噪，$t$ 次）\\[0.9ex]
  \textbf{实际上}：$\bx_0,\ \beps \xrightarrow{\ \text{式}\eqref{eq:forward-closed-reparam}\ }
                  \bx_t$ \quad（一次算出，$O(1)$）
\end{center}
训练时随机采一个 $t$、再采一个 $\beps$，直接用闭式解造出 $\bx_t$，
并以当初加进去的 $\beps$ 作为监督目标。全程只做一次加噪，与 $t$ 无关。
\end{keybox}

\subsection{噪声调度与信噪比}

$\abt$ 随 $t$ 单调递减：$t$ 越大，剩下的信号越少。刻画“信号还剩多少”的常用量是\emph{信噪比}
\begin{equation}
  \mathrm{SNR}(t)=\frac{\abt}{1-\abt},
  \label{eq:snr}
\end{equation}
它同样单调递减，从 $t=1$ 时的约 $10^4$ 一路降到 $t=1000$ 时的 $10^{-5}$ 量级。
$\{\beta_t\}$ 的取法称为噪声调度，常见两种\cite{ho2020ddpm,nichol2021improved}：

\begin{itemize}
  \item \textbf{线性调度}：$\beta_t$ 从 $\beta_1=10^{-4}$ 线性增到 $\beta_T=2\times10^{-2}$，$T=1000$。
        实现最简单，但两端变化过快、中间过于平坦；
  \item \textbf{余弦调度}：令
        $\abt=\dfrac{\cos^2\!\bigl(\frac{\pi}{2}\cdot\frac{t/T+s}{1+s}\bigr)}{\cos^2\!\bigl(\frac{\pi}{2}\cdot\frac{s}{1+s}\bigr)}$，
        $s\approx 0.008$。它让 $\abt$ 的下降更均匀，低噪声阶段保留信号更久，实践中生成质量更好。
\end{itemize}

\begin{figure}[htbp]
  \centering
  \begin{minipage}[c]{0.47\textwidth}
    \centering
    \resizebox{\linewidth}{!}{%
    \begin{tikzpicture}
      \begin{axis}[
        width=\textwidth, height=5.2cm,
        xlabel={扩散步 $t$}, ylabel={$\abt$},
        xmin=0, xmax=1000, ymin=0, ymax=1.06,
        domain=0:1000, samples=161,
        axis lines=left,
        every axis label/.style={font=\footnotesize\sffamily, text=dmNavy},
        tick label style={font=\scriptsize, text=dmGray},
        legend style={font=\scriptsize, draw=dmGray!40, fill=white,
                      at={(0.97,0.97)}, anchor=north east},
        label style={font=\scriptsize},
      ]
        \addplot[dmNavy, very thick]
          {exp(-(x*0.0001 + 0.0199*x*(x-1)/1998))};
        \addlegendentry{线性调度}
        \addplot[dmSteel, very thick, dashed]
          {(cos(deg(pi/2*(x/1000+0.008)/1.008)))^2/(cos(deg(pi/2*0.008/1.008)))^2};
        \addlegendentry{余弦调度}
      \end{axis}
    \end{tikzpicture}}
    \\[-0.4em]
    {\footnotesize\sffamily 信号保留比例 $\abt$}
  \end{minipage}\hfill
  \begin{minipage}[c]{0.47\textwidth}
    \centering
    \resizebox{\linewidth}{!}{%
    \begin{tikzpicture}
      \begin{axis}[
        width=\textwidth, height=5.2cm,
        xlabel={扩散步 $t$}, ylabel={$\mathrm{SNR}(t)$},
        xmin=1, xmax=1000, ymode=log,
        ymin=1e-5, ymax=1e5,
        domain=1:1000, samples=161,
        axis lines=left,
        every axis label/.style={font=\footnotesize\sffamily, text=dmNavy},
        tick label style={font=\scriptsize, text=dmGray},
        legend style={font=\scriptsize, draw=dmGray!40, fill=white,
                      at={(0.97,0.03)}, anchor=south east},
        label style={font=\scriptsize},
      ]
        \addplot[dmNavy, very thick]
          {exp(-(x*0.0001 + 0.0199*x*(x-1)/1998))
           /(1-exp(-(x*0.0001 + 0.0199*x*(x-1)/1998)))};
        \addlegendentry{线性调度}
        \addplot[dmSteel, very thick, dashed]
          {((cos(deg(pi/2*(x/1000+0.008)/1.008)))^2/(cos(deg(pi/2*0.008/1.008)))^2)
           /(1-((cos(deg(pi/2*(x/1000+0.008)/1.008)))^2/(cos(deg(pi/2*0.008/1.008)))^2))};
        \addlegendentry{余弦调度}
      \end{axis}
    \end{tikzpicture}}
    \\[-0.4em]
    {\footnotesize\sffamily 信噪比 $\mathrm{SNR}=\abt/(1-\abt)$}
  \end{minipage}
  \caption{两种噪声调度下 $\abt$ 与信噪比随扩散步的变化。线性调度的曲线由
    $\abt\approx\exp\bigl(-\sum_{s\le t}\beta_s\bigr)$ 给出}
  \label{fig:schedule}
\end{figure}

图~\ref{fig:schedule} 对比了两种调度。可以看到线性调度在 $t\approx 250$ 时就已经把一半以上的信号
换成了噪声，而余弦调度要到 $t\approx 500$ 才降到同一水平。这意味着余弦调度把更多的“容量”
留给了中等噪声区间，而这个区间恰恰是反向过程最难学、也最影响生成细节的部分。
